\begin{Example}[noncommutative two-tori]\label{ex:NCtori}
The noncommutative torus is a fundamental example
of a noncommutative space in both physics and mathematics. Its
original incarnation~\cite{Rieffel1981} is a nice example of a crossed
product construction, which will play a fundamental role later on in this
paper. Let $\big(C(\torus),\Z,\tau^\theta\big)$ be the commutative $C^*$-dynamical system
where $\tau^\theta$ is induced through pullback by rotations of the circle $\torus$ through
a fixed angle $\theta\in\real/\Z$:
\[
\tau^\theta_n(\mathtt{f})(z) = \mathtt{f}\big(\e^{2\pi\ii n \theta} z\big) ,
\]
for $n\in\Z$, $\mathtt{f}\in C(\torus)$ and $z\in\torus$. The resulting crossed
product
\[
\A_\theta:=C(\torus)\rtimes_{\tau^\theta}\Z
\]
is a called a \emph{rotation
algebra}, and for irrational values of $\theta$ it can be identified as a noncommutative two-torus
$\torus_\theta^2$ in the following way.

By definition, the algebra
$\A_\theta$ is the universal norm completion of the convolution
algebra $C_{\rm c}(\Z\times\torus)$, whose elements $f=\{f_n\}_{n\in\Z}$ can be regarded as
sequences (with only finitely many nonvanishing
terms) of functions $f_n\colon \torus\to\complex$. The convolution product is
given by
\[
(f\star_\theta g)_n(z) := \sum_{n'\in\Z} f_{n'}(z) g_{n-n'}\big(\e^{2\pi\ii n' \theta} z\big) ,
\]
and the $*$-algebra structure is
\[
f^\dag_n(z):=\overline{f_{-n}\big(\e^{2\pi\ii n \theta} z\big)} .
\]
Via the Fourier transformation
\[
f(z,w) := \sum_{n\in\Z} f_{n}(z) w^n
\]
for $w\in\torus$, we may regard the convolution algebra $C_{\rm
 c}(\Z\times\torus)$ as a subspace of the space of functions $C\big(\torus^2\big)$ equipped
with the star-product
\begin{gather}\label{eq:starprodNCtorus}
(f\star_\theta g)(z,w) = \sum_{n\in\Z} (f\star_\theta g)_n(z) w^n .
\end{gather}
After a further Fourier transformation
\[
f_n(z) = \sum_{m\in\Z} f_{m,n} z^m
\]
and some simple redefinitions of the Fourier series involved, the
star-product \eqref{eq:starprodNCtorus} may be written in the form
\[
(f\star_\theta g)(z,w) = \sum_{(m,n)\in\Z^2}
\bigg(\sum_{(m',n')\in\Z^2} f_{m',n'} g_{m-m',n-n'}
\e^{2\pi\ii (m-m') n' \theta}\bigg) z^m w^n .
\]
This recovers the usual commutative pointwise multiplication of functions in
$C\big(\torus^2\big)$ for $\theta=0$. For $\theta\neq0$ it realizes the irrational rotation algebra $\A_\theta$ as a
deformation of the algebra of functions $C\big(\torus^2\big)$ on a
two-dimensional torus $\torus^2$; it is equivalent to the usual strict
deformation quantization of $\torus^2$ whose star-product is a twisted
convolution product on $C\big(\torus^2\big)$.

In the language of covariant representations of the dynamical system
$(C(\torus),\Z,\tau^\theta)$, the crossed product $\A_\theta$ is the
universal $C^*$-algebra generated by two unitaries $U$ and $V$
satisfying the relation~\cite[Proposition~2.56]{Williams2007}
\begin{gather}\label{eq:C*UV}
U V=\e^{-2\pi\ii \theta} V U .
\end{gather}
A concrete representation
of $\A_\theta$ on the Hilbert space $\hil = {\rm L}^2(\torus)$ is
given by defining
\[
U(\mathtt{f})(z)=z \mathtt{f}(z) \qquad \mbox{and} \qquad V(\mathtt{f})(z) =
\mathtt{f}\big(\e^{2\pi\ii\theta} z\big) .
\]
\end{Example}

\begin{Example}[{{\bf Noncommutative $\mbf{d}$-tori}}]
\label{ex:NCtorid}
The natural higher-dimensional generalization of
Example~\ref{ex:NCtori} involves a skew-symmetric real $d{\times}d$
matrix $\Theta=(\theta_{ij})$, see~\cite{Rieffel1990}. Then the noncommutative $d$-torus
$\A_\Theta=\torus_\Theta^d$ is the universal $C^*$-algebra generated
by $d$ unitaries $U_1,\dots,U_d$ satisfying the relations
\[
U_i U_j = \e^{-2\pi\ii\theta_{ij}} U_j U_i
\]
for $i,j=1,\dots,d$. By~\cite[Lemma~1.5]{Phillips2006}, every
noncommutative torus $\torus_\Theta^d$ can be obtained as an iterated
crossed product by $\Z$ in the following way. Let
$\Theta_{|d-1}=(\theta_{ij})_{1\leq i,j\leq d-1}$, and let
$U_1,\dots,U_{d-1}$ be the standard generators of
$\A_{\Theta_{|d-1}}=\torus^{d-1}_{\Theta_{|d-1}}$. Define a
group homomorphism
\smash{$\tau^{\vec\theta}\colon \Z\to\Aut(\A_{\Theta_{|d-1}})$} by
\[
\tau_n^{\vec\theta}(U_i) = \e^{2\pi\ii n \theta_{id}} U_i
 ,
\]
for $n\in\Z$, where $\vec\theta:=(\theta_{id})\in\R^{d-1}$.
Then there is an isomorphism of $C^*$-algebras
\begin{gather}\label{eq:AThetacrossed}
\A_\Theta\simeq\A_{\Theta_{|d-1}}
\rtimes_{\tau^{\vec\theta}} \Z .
\end{gather}

In the particular case where $\Theta_{|d-1}=\mbf 0_{d-1}$, we denote the
corresponding noncommutative $d$-torus by
$\A_{\vec\theta}=\torus_{\vec\theta}^d\,$, and~\eqref{eq:AThetacrossed}
shows that it can be obtained by a crossed product of the commutative
algebra of functions on a $d{-}1$-torus by an action of~$\Z$:
\[
\A_{\vec\theta} \simeq C\big(\torus^{d-1}\big)\rtimes_{\tau^{\vec\theta}}\Z .
\]
\end{Example}

\subsection{Semi-direct products and group algebras}\label{sec:semidirect}

Most of our considerations later on will focus on spaces that can be
obtained from semi-direct products of groups. We will now explain the
relation between crossed products and semi-direct products which will
be useful for these examples.

There are two ways to think about the semi-direct product construction:
\begin{enumerate}\itemsep=0pt
\item[(1)] Let $\sfG$ be a
group with two subgroups $\sfN $ and $\sfH$ such that $\sfN $ is normal. If $\sfN \cap
\sfH=\{e\}\subset \sfG$ and every
element of $\sfG$ can be written as a product of an element of
$\sfN $ with an element of $\sfH$, then we say that $\sfG$ is a semi-direct
product of its subgroups $\sfN $ and $\sfH$ and we write $\sfG=\sfN \,\sfH$.

\item[(2)]
Let $\sfN $ and $\sfH$ be two groups together with a left action
$\varphi\colon \sfH\to\Aut(\sfN )$ of $\sfH$ on $\sfN $ by automorphisms, which we denote by
$\varphi_h(n)={}^hn$ for $h\in \sfH$ and $n\in \sfN $; in particular
${}^h(n n')={}^hn\, {}^hn'$. We write ${}^\sfH \sfN $ to
indicate that $\sfH $ acts on $\sfN $ from the left. The semi-direct product
of $\sfN $ and $\sfH $ is the group $\sfN \rtimes_\varphi \sfH $ defined to be the set
$\sfN \times \sfH $ with the product
\[
(n,h)\,(n',h') = \big(n\,{}^hn',h h'\big) .
\]
The inverse is then $(n,h)^{-1}=\big({}^{h^{-1\!}}{n^{-1}},h^{-1}\big)$.
\end{enumerate}
These two definitions are equivalent: Given subgroups $\sfN ,\sfH \subset \sfG$ as
in point (1), it follows that $\sfG\simeq \sfN \rtimes_{\rm Ad} \sfH $
where ${\rm Ad}$ is the adjoint action: ${\rm Ad}_h(n)=h n h^{-1}$. On the other
hand, every element of the group $\sfG=\sfN \rtimes_\varphi \sfH $ defined in~(2) can be written as $(n,h)=(n,e_\sfH) (e_\sfN,h)$ and the subgroups $\sfN \times\{e_\sfH\}$
and $\{e_\sfN\} \times \sfH $ intersect only in the identity of $\sfG$.
If the action of $\sfH $ on $\sfN $ is trivial,
i.e., $\varphi_h={\rm id}_\sfN $ for all $h\in \sfH $, then the semi-direct
product reduces to the direct product $\sfN \rtimes_\varphi
\sfH =\sfN \times \sfH $.

Later on we will need to consider the interplay between semi-direct
products and quotient groups, which is provided by the simple
\begin{Lemma}
\label{lem:semidirectquotient}
Let $\sfG=\sfN\rtimes_\varphi\sfH$ be a semi-direct
product, and let $\V\subset \sfN$ be a subgroup which is normal in
$\sfG$. Then the quotient group $\sfG/\V$ is the
semi-direct product $(\sfN/\V)\rtimes_{\varphi^\V}\sfH$, where
${\varphi^\V}$ is the action $\varphi$ of $\sfH$ induced on the
quotient group $\sfN/\V$.
\end{Lemma}

If $\sfN $ is a group, we write $C^*(\sfN )$ for the corresponding group
$C^*$-algebra, i.e., for the crossed product $\complex\rtimes \sfN $. If
$\sfN $ is finite, then $C^*(\sfN )=\complex[\sfN ]$ is the linear space
freely generated over $\complex$ by the group elements, made into an algebra by
linearly extending the product from $\sfN $ to $\complex[\sfN ]$; equivalently
it is the algebra of continuous functions
on $\sfN $ with the convolution product.
 Given a
left $\sfH $-action $\varphi\colon \sfH \to\Aut(\sfN )$, there is an induced action
$\varphi^*\colon \sfH \to \Aut(C^*(\sfN ))$ via pullback. For $\sfH $ and $\sfN $ finite the
vector spaces $ C^*(\sfN )\times \sfH $ and $C^*(\sfN \times \sfH )$ are canonically
isomorphic, and it is straightforward
to show that the corresponding crossed product and semi-direct product are
related by
\begin{Proposition}\label{cpsp}
If $\sfN$ and $\sfH$ are finite groups, then
\[
C^*(\sfN )\rtimes_{\varphi^*}\sfH \simeq C^*(\sfN \rtimes_\varphi \sfH ) .
\]
\end{Proposition}
\begin{proof} Note that $C^*(\sfN )\rtimes_{\varphi^*}\sfH =\complex[\sfN ]\rtimes_{\varphi^*}\sfH $ is the
 vector space of functions $f\colon \sfH \to \complex[\sfN ]$ with convolution product
\[
\big(f\star_{\complex[\sfN ]\rtimes_{\varphi^*}\sfH } f'\big)(h)=\sum_{h'\in
 \sfH }\,f(h')\star_{\complex[\sfN ]}\varphi^*_{h'}\big(f'\big(h'^{-1} h\big)\big) ,
\]
and using the convolution product in $\complex[\sfN]$ this can be
written as
\[
\big(f\star_{\complex[\sfN ]\rtimes_{\varphi^*}\sfH }
f'\big)(n,h)=\sum_{h'\in \sfH} \sum_{n'\in
 \sfN } f(n',h') f'\big({}^{h'^{-1\!}}\big(n'^{-1} n\big),h'^{-1} h\big) ,
\]
which is easily seen
 to coincide with the convolution product in $ C^*(\sfN \rtimes_\varphi \sfH )=\complex[\sfN \rtimes_\varphi \sfH ]$.
\end{proof}

A more general result holds if $\sfN $ and $\sfH $ are locally compact groups with $\varphi\colon \sfH \to \Aut (\sfN )$
a~continuous action of $\sfH $ on $\sfN $ via automorphisms
(i.e., $(h,n)\mapsto \varphi_h(n)$ is a continuous map from $\sfH \times \sfN $ to $\sfN $).
In this case the semi-direct product $\sfN \rtimes_\varphi \sfH $ is a locally compact
group (in the product topology on $\sfN \times \sfH $) with $\sfN $ a closed normal subgroup and $\sfH $ a closed
subgroup (see~\cite[Proposition~3.11]{Williams2007} with
$A=\complex$). The analogue of item~(1) above also holds in the
context of locally compact groups if $\sfG$ is $\sigma$-compact, and
$\sfN$ and $\sfH$ are closed subgroups of~$\sfG$.

The action $\beta$
defining the $C^*$-dynamical system $(C^*(\sfN ),\sfH ,\beta)$ and hence the crossed
product $C^*(\sfN )\rtimes_\beta \sfH $ is the composition
of the pullback $\varphi^*$ of the action $\varphi\colon \sfH \to \Aut(\sfN )$ with
the action $\sigma_\sfH\colon \sfH \to \real^+\subset \Aut(C^*(\sfN ))$ that enters the
definition of the Haar measure on $\sfN \rtimes_\varphi \sfH $ in terms of the Haar
measures on $\sfN $ and $\sfH $: If $\mu_\sfN $ is a (left invariant)
Haar measure on $\sfN $, then the integral
$I_h(F)=\int_\sfN F\big({}^hn\big)\, \dd\mu_\sfN (n)$ for $F\in C^*(\sfN )$
is left invariant, i.e.,
$I_h(\lambda_{n'}F)= I_h(F)$ where $(\lambda_{n'}F)(n):=F\big(n'^{-1} n\big)$
for all $n'\in \sfN $
(use invariance of the Haar measure under $n\to {}^{h^{-1\!}}n'$).
Uniqueness of the Haar measure up to a positive constant then
implies there exists a function $\sigma_\sfH\colon \sfH\to\real^+$ such that
\begin{gather}\label{sigmaI}
\sigma_\sfH(h) \int_\sfN F\big({}^hn\big) \, \dd\mu_\sfN (n)=\int_\sfN F(n)
\, \dd\mu_\sfN (n) .
\end{gather}
It is straighforward to see that
$\sigma_\sfH$ is a group homomorphism and that it is
continuous~\cite[Section~2]{Williams2007}. The Haar measure
$\mu_{\sfN\rtimes_\varphi\sfH}$ on
$\sfN \rtimes_\varphi \sfH $ is then given by
\[
\int_{\sfN\rtimes_\varphi\sfH} f(n,h) \,
\dd\mu_{\sfN\rtimes_\varphi\sfH}(n,h) := \int_\sfH \int_\sfN\,
f(n,h) \sigma_\sfH(h)^{-1} \, \dd\mu_\sfN (n) \, \dd\mu_\sfH (h) .
\]
This is trivially invariant under the left $\sfN $-action, and
it is also invariant under the left $\sfH $-action $(n,h)\mapsto
(1,h') (n,h)=\big({}^{h'\!}n, h' h\big)$, using \eqref{sigmaI} with $F\big({}^{h'\!}n\big) :=f\big({}^{h'\!}n, h' h\big)$
and recalling that~$h$ is fixed in~\eqref{sigmaI}.

\begin{Theorem}\label{cpsp2}
Let $\sfN $ and $\sfH $ be locally compact groups
 and $\varphi \colon \sfH \to \Aut(\sfN )$ a continuous action of $\sfH $ on $\sfN $. Define
$\beta\colon \sfH \to \Aut(C^*(\sfN ))$ by $(\beta_{h'}\ell)
(n)=\sigma_\sfH(h')^{-1} \ell\big({}^{h'^{-1\!}}n\big)$ for all $h'\in \sfH $ and
$\ell\in C_{\rm c}(\sfN )$. Then
\[
C^*(\sfN )\rtimes_{\beta}\sfH \simeq C^*(\sfN \rtimes_\varphi
\sfH ) .
\]
\end{Theorem}

For a full
proof of Theorem~\ref{cpsp2} that takes into account the topological
and $C^*$-algebraic aspects, see~\cite[Proposition~3.11]{Williams2007}. Here
we shall just show that under the canonical injection $C_{\rm c}(\sfN \rtimes_\varphi \sfH )\hookrightarrow
C_{\rm c}(\sfN )\rtimes_\beta \sfH $, given by $f(n,h)\mapsto f(h)$ where
$f(h)(n)=f(n,h)$, the convolution product in $C_{\rm c}(\sfN \rtimes_\varphi
\sfH )$ is mapped to the convolution product in $C_{\rm c}(\sfN )\rtimes_\beta \sfH $.
Let $f,f'\in C_{\rm c}(\sfN \rtimes_\varphi \sfH )$, then
\begin{gather}\label{eq:imagefstarf'}
\big(f\star_{C_{\rm c}(\sfN \rtimes_\varphi
 \sfH )}f'\big)(n,h)=\int_\sfH \int_\sfN f(n',h')
f'\big((n',h')^{-1} (n,h)\big) \sigma_\sfH(h')^{-1} \, \dd\mu_\sfN (n') \, \dd\mu_\sfH (h').
\end{gather}
On the other hand, for the images of $f$, $f'$ in $C_{\rm c}(\sfN
)\rtimes_\beta \sfH $ we have
\[
\big(f\star_{C_{\rm c}(\sfN )\rtimes_\beta
 \sfH }f'\big)(h)=\int_\sfH f(h')\star_{C^*(\sfN
 )}\beta_{h'}\big(f'\big(h'^{-1} h\big)\big) \, \dd\mu_\sfH (h') .
\]
Using the convolution product in $C^*(\sfN)$ this can be written as
\begin{align*}
\big(f\star_{C_{\rm c}(\sfN )\rtimes_\beta
 \sfH }f'\big)(h)(n)=\int_\sfH \int_\sfN
 f(h')(n') \beta_{h'}\big(f'\big(h'^{-1} h\big)\big)\big(n'^{-1} n\big)
\, \dd\mu_\sfN (n') \, \dd\mu_\sfH (h') ,
\end{align*}
which from the definition of the action $\beta$ is easily seen to
equal the image $(f\star_{C_{\rm c}(\sfN \rtimes_\varphi \sfH
 )}f')(h)(n)$ in $C_{\rm c}(\sfN )\rtimes_\beta
 \sfH $ of the product $(f\star_{C_{\rm c}(\sfN \rtimes_\varphi
 \sfH )}f')(n,h)$ in $C_{\rm c}(\sfN \rtimes_\varphi
\sfH )$ from~\eqref{eq:imagefstarf'}.

More generally we have~\cite[Proposition~3.11]{Williams2007}
\begin{Theorem}\label{cpsp3}
Let $(\alg,\sfN \rtimes_\varphi \sfH ,\alpha)$ be a $C^*$-dynamical
system for the semi-direct product group $\sfN \rtimes_\varphi \sfH $. Then $(\alg\rtimes_{\alpha|_\sfN }\sfN ,\sfH ,\beta)$ is a $C^*$-dynamical
system, where
\[
\beta\colon \ \sfH \longrightarrow \Aut(\alg\rtimes_{\alpha|_\sfN }\sfN ) ,
\qquad h\longmapsto \beta_h
\]
is defined by
$\big(\beta_h(f)\big)(n)=\sigma_\sfH(h)^{-1} \alpha_h\big(f\big(^{h^{-1\!}}n\big)\big)$
for all $f\in C_{\rm c}(\sfN ,\alg)\subset \alg\rtimes_{\alpha|_\sfN }\sfN $, with $\sigma_\sfH\colon
\sfH \to \real^+$ defined by~\eqref{sigmaI} and $^{h^{-1\!}}n=\varphi_{h^{-1}}(n)$.
Moreover, the canonical injection $C_{\rm c}(\sfN \rtimes_\varphi \sfH ,\alg)\hookrightarrow
C_{\rm c}\big(\sfH ,C_{\rm c}(\sfN ,\alg)\big)$
extends to a~$C^*$-algebra isomorphism
\[
\alg\rtimes_{\alpha}\big(\sfN \rtimes_\varphi
\sfH \big) \simeq \big(\alg\rtimes_{\alpha|_\sfN } \sfN
\big)\rtimes_\beta \sfH .
\]
\end{Theorem}

Theorem \ref{cpsp2} is then recovered by setting $\alg=\complex$.
